% ! TEX root = ../m3-19-20.tex

\chapter{Variabile aleatoare}

\section{Cazul discret}

\indent\indent Variabilele aleatoare sînt moduri de reprezentare a informației
în calculul probabilităților, precum matricele ne pot ajuta să \qq{codificăm}
informații geometrice sau algebrice. De obicei, ele se reprezintă asemenea
permutărilor, anume ca tablouri bidimensionale.

Începem cu definițiile fundamentale.

\begin{definition}\label{def:v-al}
  O variabile a cărei valoare se determină de un eveniment rezultat în urma
  unei experiențe se numește \emph{variabilă aleatoare}.
\end{definition}

\begin{definition}\label{def:repartitie}
  Fie $ X $ o variabilă aleatoare care poate lua valorile $ x_1, \dots, x_n $,
  cu probabilitățile $ f(x_1) $, $ \dots, f(x_n) $. Mulțimea de perechi ordonate
  $ \{ (x_i, f(x_i)) \}_i $ definește \emph{repartiția} variabilei aleatoare.

  Mai general, dacă $ X $ este o variabilă aleatoare reală (adică
  $ \mathrm{Im} f \seq \RR $), atunci funcția $ F : \RR \to \RR $ definită prin:
  \[
    F_X(x) = P(X < x), \forall x \in \RR
  \]
  se numește \emph{funcția de repartiție} a variabilei aleatoare $ X $.
\end{definition}

Putem interpreta funcția de repartiție în sensul următor: calculată în $ x $, ea
ne dă probabilitatea ca valorile variabilei aleatoare să fie pînă în $ x $.

Se poate vedea din definiție că au loc următoarele \emph{proprietăți}:
\begin{enumerate}[(1)]
\item $ \ds\lim_{x \to -\infty} F_X(x) = 0 $, iar $ \ds\lim_{x \to \infty} F_X(x) = 1 $;
\item $ F_X $ este crescătoare (deoarece acumulează) și este continuă la dreapta în
  orice punct $ x \in \RR $;
\item $ P(X = x) = F(x) - F(x - 0) $, unde $ F(x - 0) $ notează limita la stînga
  \[
    \lim_{x_0 \nearrow x} F(x_0).
  \]
\item $ P(a < X \leq b) = F(b) - F(a) $.
\end{enumerate}

Ne vor interesa două cazuri, anume acelea ale variabilelor aleatoare continue și discrete,
pe care le definim mai jos.

\begin{definition}\label{def:v-disc}
  Variabila aleatoare $ X $ se numește \emph{discretă} dacă mulțimea valorilor
  ei este o mulțime cel mult numărabile de numere reale sau complexe (în bijecție
  cu o submulțime [nestrictă] a lui $ \NN $), $ \{a_n\}_n $.
\end{definition}

În acest caz, dacă $ P(X = a_n) = p_n $, atunci condiția de normare este $ \sum p_n = 1 $,
iar funcția de repartiție $ F_X $ este o funcție în scară, cu:
\[
  F_X(x) = \sum_{a_n < x} p_n.
\]
În acest caz, putem reprezenta variabila aleatoare într-o formă matriceală:
\[
  X \sim %
  \begin{pmatrix}
    a_0 & a_1 & \dots & a_n & \dots \\
    p_0 & p_1 & \dots & p_n & \dots
  \end{pmatrix},
\]
matrice care se numește \emph{matricea de repartiție} a lui $ X $ sau \emph{distribuția} sa.

Cazul continuu este definit mai jos.

\begin{definition}\label{def:v-cont}
  Dacă $ P(X = x) = 0 $, pentru orice $ x \in \RR $, atunci variabila aleatoare $ X $ se numește
  \emph{continuă}. Echivalent, funcția ei de repartiție este o funcție reală continuă
  (v. proprietatea (3)).
\end{definition}

În multe situații, vom lucra cu variabile aleatoare discrete, reprezentate matriceal. Va fi
util să definim cîteva operații elementare între aceste obiecte.

Fie, așadar, două variabile aleatoare discrete:
\[
  X : \begin{pmatrix} x_i \\ p_i \end{pmatrix}, \quad %
  Y : \begin{pmatrix} y_j \\ q_j \end{pmatrix}
\]

Definim $ p_{ij} = P(X = x_i, Y = y_j), \forall i,j $. Atunci au loc formulele de calcul:
\begin{align*}
  p_i &= \sum_j p_{ij} \\
  q_j &= \sum_i p_{ij} \\
  1 &= \sum_{i,j} p_{ij}.
\end{align*}

Variabilele $ X $ și $ Y $ se numesc \textbf{independente} dacă $ p_{ij} = p_i \cdot q_j $.

\textbf{Suma variabilelor} are repartiția:
\[
  X + Y : \begin{pmatrix} x_i + y_j \\ p_{ij} \end{pmatrix}.
\]

Dacă $ \alpha \in \RR $ este o constantă reală, atunci variabila $ X + \alpha $ are repartiția:
\[
  X + \alpha : \begin{pmatrix} x_i + \alpha \\ p_i \end{pmatrix}.
\]

\textbf{Produsul variabilelor} $ X $ și $ Y $ are repartiția:
\[
  XY : \begin{pmatrix} x_iy_j \\ p_{ij} \end{pmatrix}
\]

\textbf{Produsul} $ \alpha X $, pentru $ \alpha \in \RR $ are repartiția:
\[
  \alpha X : \begin{pmatrix} \alpha x_i \\ p_i \end{pmatrix}.
\]

\vspace{0.3cm}

Să vedem un exemplu de calcul:

\textbf{Exemplu:} Se dau variabilele independente:
\[
  X : %
  \begin{pmatrix}
    0 & 1 & 2 \\
    0,2 & 0,4 & 0,4
  \end{pmatrix}, \quad %
  Y : %
  \begin{pmatrix}
    1 & 2 \\
    0,6 & 0,4
  \end{pmatrix}.
\]

Determinați repartițiile variabilelor $ X + Y, XY, 2 + X, X^2, 3Y $.

\emph{Soluție:} Vom prezenta repartițiile punînd valorile pentru evenimente
în ordine crescătoare. De exemplu, $ X + Y $ poate lua valorile 1, 2, 3, 4,
iar $ X^2 $ poate lua valorile 0, 1, 4.

Folosim independența și formulele de calcul pentru probabilități și obținem,
de exemplu:
\begin{align*}
  P(X + Y = 1) &= P(X = 0 \land Y = 1) \\
               &= P(X = 0) \cdot P(Y = 1) \\
               &= 0,12 \\
  P(X + Y = 2) &= P\Big((X = 0 \land Y = 2) \lor (X = 1 \land Y = 1)\Big) \\
               &= P(X = 0 \land Y = 2) + P(X = 1 \land Y = 1) \\
               &= 0,32
\end{align*}

Similar se calculează și celelalte valori și, folosind definițiile operațiilor,
obținem:
\begin{align*}
  X + Y &: \begin{pmatrix} 1 & 2 & 3 & 4 \\ 0,12 & 0,32 & 0,4 & 0,16 \end{pmatrix} \\
  XY &: \begin{pmatrix} 0 & 1 & 2 & 3 & 4 \\ 0,2 & 0,24 & 0,4 & 0,6 \end{pmatrix} \\
  2 + X &: \begin{pmatrix} 2 & 3 & 4 \\ 0,2 & 0,4 & 0,4 \end{pmatrix} \\
  X^2 &: \begin{pmatrix} 0 & 1 & 4 \\ 0,2 & 0,4 & 0,4 \end{pmatrix} \\
  3Y &: \begin{pmatrix} 3 & 6 \\ 0,6 & 0,4 \end{pmatrix}
\end{align*}

\vspace{0.3cm}


Pentru studiul variabilelor aleatoare se folosesc mai multe mărimi, pe care le introducem,
pe rînd.

\begin{definition}\label{def:medie-mom-var}
  Fie $ X $ o variabilă aleatoare discretă.
  \begin{enumerate}[(1)]
  \item Dacă seria $ \ds\sum_{n \geq 0} x_n p_n $ este absolut convergentă, atunci spunem că $ X $
    admite medie, iar numărul:
    \[
      E(X) = \sum_{n \geq 0} x_n p_n
    \]
    se numește \emph{media} lui $ X $;
  \item Pentru $ n \neq 0 $, media $ E(X^n) $ a variabilei $ X^n $ se numește \emph{momentul %
      de ordinul $ n $} al lui $ X $;
  \item Media variabilei $ (X - E(X))^2 $ se numește \emph{varianța} sau \emph{dispersia}
    lui $ X $, notată $ \mathrm{Var}(X) $ sau $ D^2(X) $.
  \end{enumerate}
\end{definition}

Funcția de medie are următoarele proprietăți esențiale:
\begin{enumerate}[(a)]
\item \textbf{Liniaritatea:} $ E(\alpha X + \beta Y) = \alpha E(X) + \beta E(Y) $, pentru
  orice $ \alpha, \beta \in \CC $;
\item \textbf{Monotonia:} Dacă $ X \leq Y $, atunci $ E(X) \leq E(Y) $;
\item Dacă $ X $ este o constantă $ c $, atunci $ E(X) = c $;
\item Dacă $ X $ și $ Y $ sînt variabile aleatoare independente, adică descriu evenimente
  independente, atunci $ E(XY) = E(X) \cdot E(Y) $.
\end{enumerate}

Funcția de dispersie are proprietățile:
\begin{enumerate}[(a)]
\item $ D^2(X) = E(X^2) - (E(X))^2 $;
\item $ D^2(X) \geq 0 $;
\item $ D^2(X) = 0 \Leftrightarrow P(X = E(X)) = 1 $, adică $ X $ este o constantă cu
  probabilitatea 1 ($ X $ descrie doar evenimentul sigur);
\item \textbf{Cebîșev:} Pentru orice $ \varepsilon > 0 $, are loc:
  \[
    P(|X - E(X)| \geq \varepsilon) < \frac{D^2(X)}{\varepsilon^2}. 
  \]
  Echivalent, se mai poate scrie:
  \[
    P(|X - E(X)| < \varepsilon) \geq 1 - \frac{D^2(X)}{\varepsilon^2}.
  \]
\item $ D^2(\alpha X) = \alpha^2 D^2(X) $, pentru orice $ \alpha \in \CC $;
\item Dacă $ X $ și $ Y $ sînt independente, atunci:
  \[
    D^2(X + Y) = D^2(X) + D^2(Y).
  \]
\end{enumerate}

Mai definim două mărimi care se asociază variabilelor aleatorii.

\begin{definition}\label{def:cov}
  Fie $ X $ și $ Y $ două variabile aleatoare. Numim \emph{covarianță} a
  variabilelor $ X $ și $ Y $, notată $ \mathrm{cov}(X, Y) $ numărul:
  \[
    \mathrm{cov}(X, Y) = E((X - E(X))(Y - E(Y))) = E(XY) - E(X)E(Y).
  \]
\end{definition}

Proprietățile esențiale ale covarianței sînt date de:
\begin{enumerate}[(a)]
\item Dacă $ X $ și $ Y $ sînt variabile aleatoare independente, atunci
  $ \mathrm{cov}(X, Y) = 0 $;
\item Dacă $ X_1, \dots, X_n $ sînt variabile aleatoare cu dispersiile
  $ D_i^2 $ și covarianțele $ \mathrm{cov}(X_i, X_j), i \neq j $, atunci are loc:
  \[
    D^2\Big( \sum_{i = 1}^n X_i \Big) = \sum_{i = 1}^n D_i^2 + %
    2 \sum_{i < j} \mathrm{cov}(X_i, X_j).
  \]
\end{enumerate}

\begin{definition}\label{def:cor}
  Dacă $ X $ și $ Y $ sînt variabile aleatoare, se numește \emph{coeficient de corelație}
  al lor, notat $ \rho(X, Y) $, numărul calculat prin:
  \[
    \rho(X, Y) = \frac{\mathrm{cov}(X, Y)}{\sqrt{D^2(X) \cdot D^2(Y)}}
  \]
\end{definition}

Se poate vedea din definiție că, atît covarianța, cît și corelația arată \qq{gradul de %
  independență} a variabilelor aleatoare $ X $ și $ Y $. Într-adevăr, dacă $ X $ și $ Y $
sînt independente, atunci avem $ \rho(X, Y) = 0 $. Afirmația reciprocă este, însă, falsă!

Să luăm un exemplu: Fie $ A = \{ a_1, a_2, a_3, a_4 \}$, cu $ P(a_i) = \dfrac{1}{4} $.

Presupunem că variabila aleatoare $ X $ realizează corespondența:
\[
  X(a_1) = 1, X(a_2) = -1, X(a_3) = 2, X(a_4) = -2.
\]
Atunci, ea are matricea de repartiție:
\[
  X \sim
  \begin{pmatrix}
    -2 & -1 & 1 & 2 \\
    \frac{1}{4} & \frac{1}{4} & \frac{1}{4} & \frac{1}{4}
  \end{pmatrix}.
\]

Presupunem că o altă variabilă aleatoare $ Y $ realizează corespondența:
\[
  Y(a_1) = 1, Y(a_2) = 1, Y(a_3) = -1, Y(a_4) = -1.
\]
Atunci matricea ei se poate reduce la:
\[
  Y \sim
  \begin{pmatrix}
    -1 & 1 \\
    \frac{1}{2} & \frac{1}{2}
  \end{pmatrix}.
\]

Produsul celor două variabile aleatoare are matricea de repartiție:
\[
  XY :
  \begin{pmatrix}
    -2 & -1 & 1 & 2 \\
    \frac{1}{4} & \frac{1}{4} & \frac{1}{4} & \frac{1}{4}
  \end{pmatrix}.
\]

Rezultă că $ E(X) = E(Y) = E(XY) $, deci $ \mathrm{cov}(X, Y) = 0 $. Însă putem
observa că $ X $ și $ Y $ nu sînt independente. Într-adevăr, avem:
\[
  P(X = 1 \land Y = 1) = P(a_1) = \frac{1}{4} \neq %
  P(X = 1) \cdot P(Y = 1) = \frac{1}{4} \cdot \frac{1}{2}.
\]

În fine, o altă noțiune de care avem nevoie este \emph{funcția generatoare}.

\begin{definition}\label{def:generatoare}
  Fie $ X $ o variabilă aleatoare discretă care ia valori naturale.

  Funcția $ G_X(t) $, definită prin:
  \[
    G_X(t) = E(t^X) = \sum_{n \geq 0} p_n t^n, 
  \]
  cu $ t \leq 1 $ și $ p_n = P(X = n) $ se numește \emph{funcția generatoare}
  a variabilei aleatoare $ X $.
\end{definition}

Funcția generatoare codifică, de fapt, în coeficienții unei serii formale,
repartiția unei variabile aleatoare. Ea are următoarele proprietăți:
\begin{enumerate}[(a)]
\item Dacă avem variabilele aleatoare independente $ X_1, \dots, X_n $,
  iar $ X = X_1 + \dots + X_n $, atunci:
  \[
    G_X(t) = G_{X_1}(t) \cdots G_{X_n}(t);
  \]
\item $ E(X) = G'_X(1) $;
\item $ D^2(X) = G_X^{\prime\prime}(1) + G'_X(1) - [G'_X(1)]^2 $;
\item Două variabile aleatoare cu aceeași funcție generatoare au aceeași repartiție.
\end{enumerate}

%%%%%%%%%%%%%%%%%%%%%%%%%%%%%%%%%%%%%%%%%%%%%%%%%%%%%%%%%%%%%%%%%%%%%%

\section{Exemple de repartiții discrete}

\indent\indent Vom regăsi acum, în contextul teoriei variabilelor aleatoare,
scheme clasice de probabilitate.



\subsection{Repartiția binomială (Bernoulli)}
\indent\indent Presupunem că se fac $ n $ experiențe independente, în fiecare din
experiențe, probabilitatea de realizare a unui eveniment $ A $ fiind constantă
și egală cu $ p $ (rezultă că probabilitatea ca $ A $ să nu se realizeze este
$ q = 1 - p $).

Numărul de realizări ale evenimentului $ A $ în cele $ n $ experiențe se constituie
ca o variabilă aleatoare $ X $, ale cărei valori posibile sînt $ k = 0, 1, \dots, n $,
în funcție de cîte din experiențe realizează $ A $. Așadar, avem $ P(X = k) = p_k $.

Cînd considerăm fiecare $ k $, cu probabilitatea sa de realizare, mulțimea perechilor
$ (k, p_k), k \in \{0, 1, \dots, n \} $ se numește \emph{repartiție binomială} sau
\emph{repartiție Bernoulli}.

Pentru acestea, valorile principalilor parametri sînt:
\begin{itemize}
\item $ P(X = k) = C_n^k p^k q^{n-k}, k \in \{0, 1, \dots, n\} $;
\item $ E(X) = np $;
\item $ D^2(X) = npq $;
\item $ G_X(t) = (pt + q)^n $.
\end{itemize}

\textbf{Observație:} Pentru $ n \to \infty, p \to 0 $, astfel încît $ np = \lambda $,
probabilitățile $ P(X = k) $ pot fi aproximate prin valorile repartiției Poisson
(vezi mai jos).

%%%%%%%%%%%%%%%%%%%%%%%%%%%%%%%%%%%%%%%%%%%%%%%%%%%%%%%%%%%%%%%%%%%%%%

\subsection{Repartiția geometrică}

\indent\indent Fie $ X $ numărul de experimente de tip Bernoulli, cu probabilitatea $ p $
de succes, care trebuie repetate pînă la apariția primului succes. Atunci:
\begin{itemize}
\item $ P(X = k) = p(1 - p)^{k - 1}, k \in \NN^\ast $;
\item $ E(X) = \dfrac{1}{p} $;
\item $ D^2(X) = \dfrac{1 - p}{p^2} $;
\item $ G_X(t) = \dfrac{pt}{1 - qt} $.
\end{itemize}

%%%%%%%%%%%%%%%%%%%%%%%%%%%%%%%%%%%%%%%%%%%%%%%%%%%%%%%%%%%%%%%%%%%%%%

\subsection{Repartiția binomial negativă}

\indent\indent Fie $ X $ numărul de experimente de tip Bernoulli cu probabilitatea
de succes $ p $ care trebuie efectuate pentru a obține $ m $ succese. Atunci:
\begin{itemize}
\item $ P(X = n) = C_{n-1}^{m - 1} p^m q^{n-m} $ (pentru $ m = 1 $, găsim repartiția
  geometrică);
\item $ E(X) = \dfrac{m}{p} $;
\item $ D^2(X) = \dfrac{mq}{p^2} $;
\item $ G_X(t) = \dfrac{p^m t^m}{(1 - qt)^m} $.
\end{itemize}

%%%%%%%%%%%%%%%%%%%%%%%%%%%%%%%%%%%%%%%%%%%%%%%%%%%%%%%%%%%%%%%%%%%%%%

\subsection{Repartiția hipergeometrică}

\indent\indent Fie $ X $ o variabilă aleatoare care reprezintă numărul de bile albe
obținute după $ n $ extrageri fără repunere dintr-o urnă ce conține $ a $ bile albe
și $ b $ bile negre. Atunci avem:
\begin{itemize}
\item $ P(X = x) = \dfrac{C_a^x C_b^{n-x}}{C_{a + b}^n}, 0 \leq x \leq n $;
\item $ E(X) = np $;
\item $ D^2(X) = xpq \cdot \dfrac{a + b - x}{a + b - 1} $, unde $ p = \dfrac{a}{a+b}$,
  iar $ q = 1 - p = \dfrac{b}{a + b} $.
\end{itemize}

%%%%%%%%%%%%%%%%%%%%%%%%%%%%%%%%%%%%%%%%%%%%%%%%%%%%%%%%%%%%%%%%%%%%%%

\subsection{Repartiția Poisson}

\indent\indent Repartiția Poisson se caracterizează printr-un parametru $ \lambda > 0 $
(v. repartiția binomială [Bernoulli] mai sus). Avem:
\begin{itemize}
\item $ P(X = n) = \dfrac{\lambda^n}{n!} \cdot e^{-\lambda}, n \in \NN $;
\item $ E(X) = \lambda $;
\item $ D^2(X) = \lambda $;
\item $ G_X(t) = e^{\lambda(t - 1)} $.
\end{itemize}

%%%%%%%%%%%%%%%%%%%%%%%%%%%%%%%%%%%%%%%%%%%%%%%%%%%%%%%%%%%%%%%%%%%%%%

\section{Cazul continuu}


\begin{definition}\label{def:va-cont}
  Variabila aleatoare $ X $ se numește \emph{absolut continuă} dacă
  există o funcție integrabilă $ f : \RR \to \RR_+ $ astfel încît funcția
  de repartiție $ F(x) $ a lui $ X $ să fie dată de:
  \[
    F(x) = \int_{-\infty}^x f(t) dt.
  \]

  În acest caz, $ f $ se numește \emph{densitatea de probabilitate} a lui $ X $.
\end{definition}

Ca în cazul discret, avem următoarele proprietăți esențiale:
\begin{enumerate}[(a)]
\item Funcția $ f $ este pozitivă și are loc condiția de normare
  \[
    \ds\int_{-\infty}^\infty f(x) dx = 1.
  \]
\item În orice punct de continuitate a lui $ f $, avem $ F'(x) = f(x) $;
\item $ P(a \leq X \leq b) = \ds\int_a^b f(x) dx $;
\item $ E(X) = \ds\int_{-\infty}^\infty xf(x) dx $;
\item Dacă $ X $ are densitatea $ X $, iar $ Y = aX + b $ este o altă variabilă
  aleatoare, cu $ a, b \in \RR, a \neq 0 $, atunci $ Y $ are densitatea:
  \[
    f_Y(x) = \frac{1}{|a|} f \Big( \frac{x - b}{a} \Big).
  \]
\end{enumerate}

\begin{definition}\label{def:fc-car}
  Pentru orice variabilă aleatoare $ X $, funcția:
  \[
    \varphi_X(t) = E(e^{itX})
  \]
  se numește \emph{funcția caracteristică} a lui $ X $.
\end{definition}

Funcția caracteristică joacă un rol similar funcției generatoare din cazul discret.
Astfel, ea are următoarele proprietăți:
\begin{enumerate}[(a)]
\item Dacă $ X_1, \dots, X_n $ sînt variabile aleatoare independente, iar
  $ X = X_1 + \dots + X_n  $, atunci:
  \[
    \varphi_X(t) = \prod_{k = 1}^n \varphi_{X_k}(t);
  \]
\item Dacă $ X $ admite moment de ordinul $ n $, atunci:
  \[
    E(X^k) = \frac{\varphi_X^{(k)}(0)}{i^k}, \quad 1 \leq k \leq n;
  \]
\item Două variabile aleatoare cu aceeași funcție caracteristică au aceeași repartiție;
\item Dacă $ X $ este discretă, atunci obținem:
  \[
    \varphi_X(t) = \sum_n e^{ita_n} p_n;
  \]
\item Dacă $ X $ are densitatea de repartiție $ f(x) $, atunci:
  \[
    \varphi_X(t) = \int_{-\infty}^\infty e^{itx} f(x) dx,
  \]
  iar în punctele de continuitate ale lui $ f $, are loc relația:
  \[
    f(x) = \frac{1}{2\pi} \int_{-\infty}^\infty e^{-itx} \varphi_X(t) dt.
  \]
\end{enumerate}

În continuare, prezentăm cele mai cunoscute repartiții absolut continue.

%%%%%%%%%%%%%%%%%%%%%%%%%%%%%%%%%%%%%%%%%%%%%%%%%%%%%%%%%%%%%%%%%%%%%%

\section{Repartiții absolut continue}

\subsection{Repartiția normală (gaussiană)}

\indent\indent Dată media $ m $ și abaterea pătratică $ \sigma $, repartiția
variabilei aleatoare $ X $ are densitatea:
\[
  f(x) = \frac{1}{\sigma \sqrt{2\pi}} \cdot e^{-\frac{(x - m)^2}{2\sigma^2}}.
\]

În acest caz, se mai notează $ X \sim N(m, \sigma) $.

Pentru cazul particular $ m = 0 $ și $ \sigma = 1 $, $ X $ se numește
\emph{variabilă aleatoare standard}.

Funcția caracteristică este dată de $ \varphi_X(t) = e^{imt - \frac{\sigma^2 t^2}{2}}. $

Pentru cazul variabilelor normale standard, funcția de repartiție are o formă
particulară și se numește \emph{funcția lui Laplace}:
\[
  \Phi(x) = \frac{1}{\sqrt{2\pi}} \int_{-\infty}^x e^{-\frac{t^2}{2}} dt.
\]

\textbf{Observație:} Valorile funcției Laplace sînt tabelate și, de obicei,
se dau în probleme. \\

În general, dacă avem o variabilă aleatoare normală de tip $ N(m, \sigma) $, atunci
funcția ei de repartiție $ F $ se poate obține din funcția lui Laplace în forma:
\[
  F(x) = \Phi \Big( \frac{x - m}{\sigma} \Big).
\]

Așadar, avem:
\begin{equation}
  \label{eq:p-phi-lap}
  P(a \leq X \leq b) = \Phi \Big( \frac{b - m}{\sigma} \Big) - \Phi \Big( \frac{a - m}{\sigma} \Big),
\end{equation}
de unde putem obține mai departe:
\[
  P(|X - m| < \varepsilon \sigma) = \Phi(\varepsilon) - \Phi(-\varepsilon) = %
  2 \Phi(\varepsilon) - 1.
\]


%%%%%%%%%%%%%%%%%%%%%%%%%%%%%%%%%%%%%%%%%%%%%%%%%%%%%%%%%%%%%%%%%%%%%%

\subsection{Repartiția uniformă}

\indent\indent Pentru un interval $ [a, b] $, densitatea de repartiție este:
\[
  f(x) = \begin{cases}
    \dfrac{1}{b - a}, & a \leq x \leq b \\
    0, & \text{ în rest}
  \end{cases}.
\]

În acest caz, avem media $ E(X) = \dfrac{a + b}{2} $ și
$ D^2(X) = \dfrac{(b - a)^2}{12} $.


%%%%%%%%%%%%%%%%%%%%%%%%%%%%%%%%%%%%%%%%%%%%%%%%%%%%%%%%%%%%%%%%%%%%%%

\subsection{Repartiția exponențială}

Dat un parametru $ \lambda > 0 $, repartiția exponențială are densitatea:
\[
  f(x) = \begin{cases}
    \lambda e^{-\lambda x}, & x > 0 \\
    0, & x \leq 0
  \end{cases}
  .
\]

Pentru aceasta, avem:
\begin{itemize}
\item $ E(X) = \dfrac{1}{\lambda} $;
\item $ D^2(X) = \dfrac{1}{\lambda^2} $;
\item $ \varphi_X(t) = \Big( 1 - \dfrac{it}{\lambda} \Big)^{-1} $.
\end{itemize}

%%%%%%%%%%%%%%%%%%%%%%%%%%%%%%%%%%%%%%%%%%%%%%%%%%%%%%%%%%%%%%%%%%%%%%

\subsection{Repartiția Gamma}

\indent\indent Dați parametrii $ \lambda, p > 0 $, densitatea variabilei
aleatoare de repartiție Gamma este:
\[
  f(x) = \begin{cases}
    \dfrac{\lambda^p}{\Gamma(p)} x^{p-1} e^{-\lambda x}, & x > 0 \\
    0, & x \leq 0
  \end{cases},
\]
unde $ \Gamma $ este funcția lui Euler. În acest caz, avem:
\begin{itemize}
\item $ E(X) = \dfrac{p}{\lambda} $;
\item $ D^2(X) = \dfrac{p}{\lambda^2} $;
\item $ \varphi_X(t) = \Big( 1 - \dfrac{it}{\lambda} \Big)^{-p} $.
\end{itemize}

Din repartiția Gamma, putem obține alte cîteva cazuri particulare:
\begin{itemize}
\item pentru $ p = 1 $, se obține repartiția exponențială;
\item pentru $ \lambda = \dfrac{1}{2}, p = \dfrac{n}{2} $, se obține
  repartiția numită \emph{hi pătrat}, cu $ n $ grade de libertate,
  notată $ \chi^2(n) $.
\end{itemize}

%%%%%%%%%%%%%%%%%%%%%%%%%%%%%%%%%%%%%%%%%%%%%%%%%%%%%%%%%%%%%%%%%%%%%%

\subsection{Repartiția Cauchy}

\indent\indent Aceasta are densitatea $ f(x) = \dfrac{1}{\pi(1 + x^2)} $.

În acest caz, $ X $ nu admite valoare medie.

%%%%%%%%%%%%%%%%%%%%%%%%%%%%%%%%%%%%%%%%%%%%%%%%%%%%%%%%%%%%%%%%%%%%%%

\section{Exerciții}

\indent\indent 1. Doi parteneri egal cotați boxează 12 runde. Probabilitatea
ca oricare dintre ei să cîștige o rundă este 50\%. Să se calculeze valoarea
medie, dispersia și abaterea medie pătratică a variabilei aleatoare care
reprezintă numărul de runde cîștigate de unul din parteneri.

\emph{Soluție:} Variabila aleatoare are o distribuție binomială.
Astfel, avem:
\begin{itemize}
\item $ P(X = k) = C_{12}^k \dfrac{1}{2^k} \cdot \dfrac{1}{2^{12-k}} $;
\item $ E(X) = np = 12 \cdot \dfrac{1}{2} $;
\item $ D^2(X) = npq = 12 \cdot \dfrac{1}{2} \cdot \frac{1}{2} = 3 $;
\item $ D(X) = \sqrt{D^2(X)} = \sqrt{3} $.
\end{itemize}

\vspace{0.3cm}

2. La o agenție de turism s-a constatat că 5\% dintre persoanele care au făcut
rezervări renunță. Se presupune că s-au făcut 100 de rezervări pentru un hotel cu
95 locuri. Care este probabilitatea ca toate persoanele ce se prezintă la hotel
să aibă loc?

\emph{Soluție:} Fie $ X $ variabila aleatoare corespunzătoare numărului de persoane
care se prezintă la hotel. $ X $ urmează o repartiție binomială, cu $ n = 100 $ și
$ p = 0,95 $. Atunci avem:
\[
  P(X \leq 95) = 1 - P(X > 95) = %
  1 - \sum_{k = 96}^{100} C_{100}^k (0,05)^{100-k} \cdot (0,95)^k.
\]

\vspace{0.3cm}

3. La examenul de matematică, probabilitatea ca o teză să fie notată cu notă de
trecere este 75\%. Se aleg la întîmplare 10 lucrări. Fie $ X $ variabila aleatoare
ce reprezintă numărul tezelor ce vor fi notate cu notă de trecere. Calculați:
\begin{enumerate}[(a)]
\item repartiția lui $ X $ ;
\item $ E(X), D^2(X) $;
\item $ P(X \geq 5), P(7 \leq X \leq 10 / X \geq 8), P(X = 10) $;
\item funcția caracteristică a lui $ X $.
\end{enumerate}

\emph{Soluție:} (a) Avem o variabilă aleatoare cu repartiție binomială și $ n = 10 $,
$ p = 0,75 $. Atunci $ X $ ia valorile $ x \in \{0, 1, \dots, 10 \}$ cu probabilitățile
corespunzătoare distribuției binomiale, deci:
\[
  p(x) = C_{10}^x (0,75)^x \cdot (0,25)^{10-x}.
\]

(b) Conform formulelor de calcul, avem:
\begin{align*}
  E(X) &= np = 10 \cdot 0,75 = 7,5 \\
  D^2(X) &= npq = 10 \cdot 0,75 \cdot 0,25 = 1,875.
\end{align*}

(c)
\[
  P(X \geq 5) = \ds\sum_{x = 5}^{10} C_{10}^x (0,75)^x \cdot (0,25)^{10-x}
\]

Pentru $ P(7 \leq X \leq 10 / X \geq 8) $ avem un \emph{eveniment condiționat}.
Dacă $ A $ și $ B $ sînt două evenimente, amintim că probabilitatea $ P(A/B) $,
adică probabilitatea evenimentului $ A $ condiționat de evenimentul prealabil $ B $
este:
\[
  P(A / B) = \frac{P(A \cap B)}{P(B)}.
\]

Atunci, obținem:

\begin{align*}
  P(7 \leq X \leq 10 / X \geq 8) &= \dfrac{P(8 \leq X \leq 10)}{P(X \geq 8)} \\
              &= \frac{\sum_{x = 8}^{10} C_{10}^x (0,75)^x (0,25)^{10-x}}{%
                \sum_{x = 8}^{10} C_{10}^x (0,75)^x (0,25)^{10-x}} \\
              &= 1
\end{align*}

În fine:
\[
  P(X = 10) = C_{10}^{10} (0,75)^x (0,25)^0.
\]

(d) Calculăm succesiv:
\begin{align*}
  \varphi_X(t) &= E(e^{itX}) \\
               &= \sum_{x = 0}^{10} e^{itx} C_{10}^x (0,75)^x (0,25)^{10-x} \\
               &= (0,75 \cdot e^{it} + 0,25)^{10}.
\end{align*}

\vspace{0.3cm}

4. Presupunem că la 100 de convorbiri telefonice au loc 1000 de bruiaje.
Care este probabilitatea de a avea o convorbire fără bruiaje? Dar una
cu cel puțin 2 bruiaje?

Variabila aleatoare ce reprezintă numărul de bruiaje se consideră a fi
repartizată Poisson, cu $ \lambda = E(X) $.

\emph{Soluție:} Avem $ E(X) = \dfrac{1000}{100} = 10 = \lambda $, care este
și parametrul distribuției Poisson. Atunci:
\begin{align*}
  P(X = 0) &= e^{-10} \cdot \frac{10^0}{0!} = e^{-10} \\
  P(X \geq 2) &= 1 - P(X = 0) - P(X = 1) \\
           &= 1 - e^{-10} - 10 \cdot e^{-10}.
\end{align*}

\vspace{0.3cm}

5. Într-o mină au loc în medie 2 accidente pe săptămînă, distribuite după legea
Poisson. Să se calculeze probabilitatea de a avea cel mult 2 accidente:
\begin{enumerate}[(a)]
\item într-o săptămînă;
\item în 2 săptămîni;
\item în fiecare săptămînă dintr-un interval de 2 săptămîni.
\end{enumerate}

\emph{Soluție:} (a) Fie $ X $ variabila aleatoare ce desemnează numărul de accidente
dintr-o săptămînă. Distribuția este de tip Poisson, cu $ \lambda = E(X) = 2 $.
Atunci:
\[
  P(X \leq 2) = e^{-2} \Big( 1 + \frac{2}{1!} + \dfrac{4}{2!} \Big) = 5e^{-2}.
\]

(b) Fie $ Y $ variabila aleatoare ce desemnează numărul de accidente în 2 săptămîni.
Distribuția este de asemenea de tip Poisson, cu $ \lambda = 2 + 2 = 4 $. Atunci:
\[
  P(Y \leq 2) = e^{-5} \Big( 1 + \frac{4}{1!} + \frac{16}{2!} \Big) = 13e^{-4}.
\]

(c) Probabilitatea cerută este:
\[
  P(X \leq 2) \cdot P(X \leq 2) = 25e^{-4}.
\]

6. La un anumit cîntar erorile de măsurare sînt normal distribuite cu $ m = 0 $ și
$ \sigma = 0,1 g $. Dacă se cîntărește un obiect la acest aparat, care este
probabilitatea ca eroarea de măsurare să fie mai mică decît 15g?

\emph{Soluție:} Fie $ X $ variabila aleatoare care reprezintă eroarea de măsurare
dată cînd un obiect este cîntărit. Această variabilă este continuă și distribuită
normal cu parametrii dați. Notăm acest lucru cu $ X \sim N(0; 0,1) $.
Vrem să calculăm probabilitatea $ P(-0,15 \leq X \leq 0,15) $.

Avem succesiv (cf.~\eqref{eq:p-phi-lap}):
\begin{align*}
  P(-0,15 \leq X \leq 0,15) &= P\Big( \frac{-0,15 - 0}{0,1} \leq X \leq \frac{-0,15 + 0}{0,1} \Big) \\
                            &= P(-1,5 \leq X \leq 1,5) \\
                            &= \Phi(1,5) - \Phi(-1,5) \\
                            &= \Phi(1,5) - 1 + \Phi(1,5) \\
                            &= 0,8664.
\end{align*}

\textbf{Observație:} Cum exercițiile din acest seminar sînt rezolvate integral,
pentru temă, puteți alege exerciții nerezolvate fie din notițele doamnei profesoare
Costache (paginile 102--105) sau din cartea de la Universitatea din Iași (paginile 82--87).
Chiar dacă exercițiile au răspunsuri, ca și pînă acum, soluțiile trebuie însoțite
de un videoclip care să le explice.


%%% Local Variables:
%%% mode: latex
%%% TeX-master: "../m3-19-20"
%%% End:
